\documentclass[11pt,a4paper]{article}
\usepackage[utf8]{inputenc}
\usepackage[T1]{fontenc}
\usepackage{lmodern}
\usepackage[a4paper,margin=22mm]{geometry}
\usepackage{graphicx}
\usepackage{booktabs}
\usepackage{amsmath,amssymb}
\usepackage{xcolor}
\usepackage{caption}
\usepackage{hyperref}
\usepackage{tabularx}
\usepackage{array}
\hypersetup{colorlinks=true,linkcolor=blue!50!black,urlcolor=blue!50!black}

\title{Gen-3 Milestone-1 Track-Extrapolation Models:\\
       Independent Re-Evaluation and Allen Compatibility Audit}
\author{TrackExtrapolation / gen-3 \\ \small auto-generated from gen3\_model\_analysis.ipynb}
\date{2026-05-08}

\begin{document}
\maketitle

\begin{abstract}
We re-evaluate the five trained Gen-3 Milestone-1 (M1) track-extrapolation
models on a held-out test split using a single uniform analysis pipeline,
independently of the per-run training scripts. We confirm the headline
result that the \texttt{nrk4\_tiny\_1step\_v1} model achieves
$\langle|\Delta x|\rangle \approx 113\,\mu\mathrm{m}$ at
$\sim$4\,997 trainable parameters --- two orders of magnitude better than
either pure-MLP control --- and report a full Allen Stage-1 ABI compatibility
audit. Of the five candidates, only \texttt{nrk4\_tiny\_1step\_v1} satisfies
all four hard ABI gates (A2 bias, A3 size, A5 outputs, A7 inputs) at
fp32; two more become compliant at fp16. None of the five reaches the
Stage-1 \emph{physics} tolerance of $|\Delta x|<12\,\mu\mathrm{m}$ at
the VELO exit, so M1 is architecturally green-lit but performance red-lit.
We document an eight-action backlog (M1.5) targeting the residual
heavy-tail, NRK4 backward asymmetry, PINN $q/p$ bias, and the move from
200\,k to the full 10\,M training corpus.
\end{abstract}

\section{Models under review}
Five checkpoints from \texttt{experiments/gen\_3/trained\_models/} are
loaded directly: two pure MLPs (control), two Neural-RK4 variants and one
PINN-v2. All models share a common 7-input / 5-output Allen V3 contract:
$(x, y, t_x, t_y, q/p, z_0, dz)\,\to\,(x_f, y_f, t_{x,f}, t_{y,f}, q/p)$.

\section{Methods}
The full 10\,M-track Gen-3 corpus is loaded once. Each model is evaluated
on its persisted \texttt{test\_indices.npy} subset (20\,k tracks). All
reported metrics --- mean / median / 95th-percentile $|\Delta x|$, station
gates, the SciFi-T3 $q/p$ correlation $\rho$, and the backward/forward
ratio --- are computed by a single pipeline in the analysis notebook,
independent of per-run training scripts. Throughput is measured on a 24-core
EPYC CPU.

\section{Results}

\subsection{Headline residuals}
\begin{table}[h!]
\centering
\caption{Global residual statistics on the 20\,k-track test split.
$|\Delta x|$ in $\mu$m, throughput in tracks/s.}
\small
\begin{tabular}{lllrrrrrr}
\toprule
Model & Type & Params & Mean & Median & p95 & RMS pos. & Throughput \\
\midrule
mlp\_small\_v1 & mlp & 18,052 & 18408.0 & 11818.6 & 56952.7 & 36348.3 & 344,497 \\
mlp\_medium\_v1 & mlp & 51,076 & 5569.3 & 3684.8 & 17047.3 & 10969.0 & 414,203 \\
nrk4\_tiny\_1step\_v1 & neural\_rk4 & 4,997 & 112.9 & 76.9 & 212.2 & 658.4 & 115,636 \\
nrk4\_small\_1step\_v1 & neural\_rk4 & 18,181 & 210.0 & 160.6 & 470.5 & 688.9 & 118,810 \\
pinn\_v2\_small\_v1 & pinn\_v2 & 10,372 & 116.7 & 11.7 & 454.1 & 646.3 & 745,462 \\
\bottomrule
\end{tabular}
\end{table}

\subsection{Per-station Stage-1 gates}
\begin{table}[h!]
\centering
\caption{Stage-1 station-binned $\langle|\Delta x|\rangle$ ($\mu$m), the
SciFi-T3 $q/p$-correlation, and the bwd/fwd ratio.
Allen Stage-1 thresholds: VELO 12\,$\mu$m, UT 50\,$\mu$m,
SciFi $|\rho|<0.1$, ratio $\in[0.5,2.0]$.}
\small
\begin{tabular}{lrrrrr}
\toprule
Model & VELO exit & UT entry & SciFi T3 & $\rho_{T3}$ & bwd/fwd \\
\midrule
mlp\_small\_v1 & 20790.7 & 19181.8 & 17296.1 & +0.029 & 0.92 \\
mlp\_medium\_v1 & 6981.6 & 5867.8 & 4940.4 & -0.017 & 1.03 \\
nrk4\_tiny\_1step\_v1 & 112.6 & 115.5 & 120.3 & -0.025 & 1.44 \\
nrk4\_small\_1step\_v1 & 199.2 & 207.1 & 217.7 & -0.032 & 1.64 \\
pinn\_v2\_small\_v1 & 165.0 & 124.1 & 124.2 & +0.217 & 0.90 \\
\bottomrule
\end{tabular}
\end{table}

\subsection{Allen ABI audit}
\begin{table}[h!]
\centering
\caption{Allen Stage-1 ABI audit. Column \emph{overall} requires A2, A3 fp32,
A5, A7 simultaneously.}
\small
\begin{tabular}{lrrrcccccc}
\toprule
Model & kB fp32 & kB fp16 & $\rho_{T3}$ & A2 & A3$_{32}$ & A3$_{16}$ & A5 & A7 & overall \\
\midrule
mlp\_small\_v1 & 70.5 & 35.3 & +0.029 & \checkmark & $\times$ & \checkmark & \checkmark & \checkmark & $\times$ \\
mlp\_medium\_v1 & 199.5 & 99.8 & -0.017 & \checkmark & $\times$ & $\times$ & \checkmark & \checkmark & $\times$ \\
nrk4\_tiny\_1step\_v1 & 19.5 & 9.8 & -0.025 & \checkmark & \checkmark & \checkmark & \checkmark & \checkmark & \checkmark \\
nrk4\_small\_1step\_v1 & 71.0 & 35.5 & -0.032 & \checkmark & $\times$ & \checkmark & \checkmark & \checkmark & $\times$ \\
pinn\_v2\_small\_v1 & 40.5 & 20.3 & +0.217 & $\times$ & \checkmark & \checkmark & \checkmark & \checkmark & $\times$ \\
\bottomrule
\end{tabular}
\end{table}

\subsection{Figures}

\begin{figure}[h!]
\centering
\includegraphics[width=0.85\linewidth]{figures/gen3_m1/fig1_residual_hist.png}
\caption{Per-track $|\Delta x|$ distribution on the test split (log-x).
Reference lines mark the VELO 12\,$\mu$m gate and the gen-2 reference
of 125\,$\mu$m.}
\end{figure}

\begin{figure}[h!]
\centering
\includegraphics[width=0.85\linewidth]{figures/gen3_m1/fig2_zf_bins.png}
\caption{Mean $|\Delta x|$ vs propagation endpoint $z_f$ in eight bins.
Both NRK4 variants and the PINN remain in the 100--200\,$\mu$m band across
the full LHCb $z$ range; both MLPs sit two to three orders of magnitude
above. The shaded band is the Stage-1 VELO gate.}
\end{figure}

\begin{figure}[h!]
\centering
\includegraphics[width=\linewidth]{figures/gen3_m1/fig3_dx_vs_qop_scifi.png}
\caption{$\Delta x$ vs $q/p$ in the SciFi T3 window (8500 $\le z_f<$ 10000\,mm).
The $\rho$ value annotated on each panel is the Allen A2 metric;
\texttt{pinn\_v2\_small\_v1} fails the $|\rho|<0.1$ tolerance.}
\end{figure}

\section{Discussion}

\paragraph{Architecture beats capacity.}
\texttt{nrk4\_tiny\_1step\_v1} (4\,997 weights, 19.5\,kB fp32) outperforms
\texttt{mlp\_medium\_v1} (51\,076 weights, 199.5\,kB fp32) by a factor
$\sim$\,49 on mean $|\Delta x|$. The physics-informed forward operator
(one RK4 step driven by the analytic LHCb field, accessed through a
non-trainable buffer) is doing the heavy lifting; the corrector net is
small, fast, and easily Allen-compliant.

\paragraph{Two MLPs are not deployment candidates.}
Mean $|\Delta x|$ of 5.6\,mm and 18.4\,mm is two orders above
the gen-2 baseline. Pure MLPs serve only as a control group documenting
that field information is essential.

\paragraph{Allen footprint discriminates the field.}
At fp32 only \texttt{nrk4\_tiny} satisfies A3 ($\le$\,64\,kB).
\texttt{mlp\_small\_v1} and \texttt{nrk4\_small\_1step\_v1} fit only at
fp16; they are deployment candidates contingent on a $\mu$m-level fp16
round-trip test.

\paragraph{PINN\_v2 fails the bias gate, not the size gate.}
Despite an excellent median residual of $\sim$\,12\,$\mu$m and a 40-kB
weight footprint, the SciFi-T3 $\rho(\Delta x, q/p) = +0.217$ exceeds the
Allen A2 tolerance. The PINN audit attributes this to (i)~dimensional
inconsistency in the PDE residual scaling, (ii)~a tautological
initial-condition term, and (iii)~a corrector head decoupled from the
magnet kick.

\paragraph{Heavy-tailed residuals everywhere.}
Across all five models the mean / median ratio is 1.4--10. Squared-MSE is
therefore unsuitable as the model-selection metric on signed-$dz$ data.
M1.5 will adopt a Huber + log-cosh blend (action \textbf{M}).

\section{Next steps -- M1.5 backlog}
\begin{description}
\item[\textbf{M.}] Replace MSE with a median-aware loss
(Huber + log-cosh). Acceptance: mean/median ratio $<2$ on val.
\item[\textbf{N.}] Move corrector net outside RK4 \texttt{\_rhs}; gate by
FSAL $\tanh(\alpha\|k_4-k_3\|\,|h|)$. Acceptance: bwd/fwd $\in[0.85,1.15]$.
\item[\textbf{O.}] PINN-v2 head refactor: dimensional consistency in PDE
residual; remove tautological IC; couple corrector to magnet kick.
Acceptance: $|\rho_{T3}|<0.05$.
\item[\textbf{P.}] Train on full 10\,M tracks for $\ge200$ epochs with cosine
schedule. Acceptance: mean $|\Delta x|<25\,\mu$m at SciFi T3.
\item[\textbf{Q.}] Export Allen \texttt{\_\_constant\_\_} weight blob for
\texttt{nrk4\_tiny}; verify fp16 round-trip $<1\,\mu$m.
\item[\textbf{R.}] Add Allen kernel skeleton (7\,$\to$\,5 entry kernel,
B-grid texture, FSAL switch). Acceptance: compiles + Allen unit-test runs.
\item[\textbf{S.}] Train on $z_0\in[0,8000]$, test on $[8000,14000]$.
Acceptance: mean $|\Delta x|$ does not exceed in-distribution by $>2\times$.
\item[\textbf{T.}] Re-run the full set with fixed seeds + MLflow tracking;
log all sweeps including failed runs.
\end{description}

\section*{Reproducibility}
This report is regenerated by re-running
\texttt{experiments/gen\_3/notebooks/gen3\_model\_analysis.ipynb} end-to-end.
Inputs: \texttt{trained\_models/<run>/best\_model.pt},
\texttt{normalization.json}, \texttt{config.json},
\texttt{test\_indices.npy}, \texttt{data/train\_10M\_gen3.npz}. Generated
artifacts: \texttt{notebooks/figures\_gen3\_m1/\{headline\_metrics,
stage1\_gates, allen\_audit\}.csv} and three PNGs.

\end{document}
